\hypertarget{correlation}{%
\section{Replacing prime and roughness tests}\label{correlation}}

We replace prime and roughness indicators by their complete-cell
averages. The weights \(A_\rho\), \(x=e^L\), and \(X_\rho\) are those of
\protect\hyperlink{proof}{the main proof}. All constants are fixed
before taking \(L\to\infty\). In particular, a constant attached to a
finite partition is not allowed to vary with \(L\).

The estimates concern the particular prime and roughness discrepancies
defined below. The other coefficient in a bilinear sum may be an
arbitrary bounded complex function. This is the property needed when
replacing the second of two factors.

\Needspace{6\baselineskip}

\hypertarget{correlation-complete-logarithmic-cells}{%
\subsection{Complete logarithmic
cells}\label{correlation-complete-logarithmic-cells}}

Put \(W=\exp(\sqrt L)\) and \(\Delta=L^{-S}\), where \(S>0\) will be
chosen. For \(j\ge0\) define

\[
C_j=(e^{j\Delta},e^{(j+1)\Delta}]\cap\mathbb N.
\]

Every integer \(n\ge2\) belongs to the unique cell with
\(j=\lceil\log n/\Delta\rceil-1\). Define the roughness indicator

\[
R_y(n)=1_{n\ge2,\ P^-(n)>y}.
\]

For either \(I(n)=R_y(n)\), with \(y\ge W\), or
\(I(n)=1_{n\text{ prime}}\), define its approximation by

\[
B_I(n)=R_W(n)\,
 \frac{\sum_{m\in C_j}I(m)}{\sum_{m\in C_j}R_W(m)}
 \quad(n\in C_j),\qquad \alpha_I(n)=I(n)-B_I(n). \tag{C1}
\]

A zero denominator is assigned the value zero. On a cell where \(I\) is
supported on the \(W\)-rough integers, a zero denominator also forces
the numerator to vanish. Consequently

\[
\sum_{n\in C_j}\alpha_I(n)=0. \tag{C2}
\]

Restricting a coefficient to a factor interval means multiplying it by
that interval's indicator. The ratio in (C1) continues to use the whole
cell, including when the interval cuts one or both endpoint cells. This
convention makes restriction to a smaller interval an exact identity.

For the principal character, (C2) handles complete cells. On each such
cell, subtract the value of \(n^{-1+it}\) at one endpoint. The
derivative bound

\[
\left|\frac{d}{dn}n^{-1+it}\right|\le\frac{1+|t|}{n^2}
\]

bounds the variation. At most two endpoint cells remain. Since
\(q\le W\), the principal character is one on every supported integer.
The resulting bound is

\[
\left(8(e^\Delta-1)+\frac2a\right)(2+|t|).
\]

For a nonprincipal character, write each \(y\)-rough integer uniquely as
\(n=pd\), where \(p=P^-(n)\). For fixed \(d\), the prime \(p\) lies in
the interval obtained from \(a\le pd\le b\), \(p>y\), and
\(p\le P^-(d)\); when \(d=1\) the last condition is omitted. This is a
single clipped interval, even when \(p\) divides \(d\). Siegel--Walfisz
and partial summation bound its prime character sum with weight
\(p^{-1+it}\) by \(K(2+|t|)(\log y)^{-D'}\) for any requested \(D'\),
uniformly when the modulus is at most a fixed power of \(\log y\). The
outer cofactor costs only

\[
\sum_{d\le b}\frac1d\le1+\log b.
\]

Because \(\log y\ge\sqrt L\), \(b\le e^{2L}\), and \(q,|t|\le L^B\),
choosing \(D'\) sufficiently large gives any fixed inverse power of
\(L\) for the rough character sum. On each cell slice, the proxy is a
scalar in \([0,1]\) times the \(W\)-rough indicator, so the same
estimate applies. There are at most \(2L^{S+1}+2\) slices. Requesting
saving \(L^{-(B+S+4)}\) for the rough sums makes the total
test-minus-proxy sum at most \(L^{-B}\).

Combining the two character cases gives the following precise statement:
for fixed \(B>0\), \(S>2B\), and \(\beta>0\), eventually, uniformly for

\[
2\le a\le b\le2a,\quad e^{\beta L}\le a,\quad b\le e^{2L},\quad
1\le q\le L^B,\quad |t|\le L^B,\quad y\ge W,
\]

and every Dirichlet character \(\chi\) modulo \(q\),

\[
\left|\sum_{a\le n\le b}\alpha_{R_y}(n)\chi(n)n^{-1+it}\right|
\le L^{-B}. \tag{C3}
\]

For the prime discrepancy, use \(y=2\sqrt a\): every complete cell
meeting \([a,b]\) lies in a range where \(y\)-roughness is equivalent to
primality. The two proxies therefore agree on the complete cells,
including the two cells cut by the interval endpoints.

\Needspace{6\baselineskip}

\hypertarget{correlation-the-bilinear-estimate-used-in-the-proof}{%
\subsection{The bilinear estimate used in the
proof}\label{correlation-the-bilinear-estimate-used-in-the-proof}}

For finite sets of positive integers \(K,J\), set

\[
\mathcal T_\rho(K,J;\alpha,\beta)
 =\sum_{0\le u\le Y}w_\rho(u)
   \sum_{\substack{m\in K,\ n\in J\\mn=u+2}}\alpha(m)\beta(n),
\qquad Y=\lfloor2x\rfloor. \tag{C4}
\]

All pairs are ordered, and \(m=n\) is retained. In applications
\(\alpha\) is a restricted discrepancy (C1), while \(|\beta|\le1\). The
\protect\hyperlink{bilinear-the-statement-on-one-factor-rectangle}{bilinear
chapter} states the estimate with all factor ranges and quantifiers: on
an admissible dyadic rectangle it bounds (C4) by \(CxL^{-D}\) for any
prescribed \(D\), with the cell mesh and \(C\) chosen before the
interval endpoints and the bounded companion coefficient.

Here is the parameter choice used to apply it. Fix \(0<b_1<1/100\) and
\(0<\kappa<1/50\), and put

\[
\beta_P=b_1/12,\qquad \beta=b_1/3,\qquad b_*=b_1/3.
\]

Reserve one ordinary prime slot in a fixed band \(Q_{j_Q}\) so small
that

\[
\frac65\sigma2^{-j_Q}\le\frac{\beta_P}{4}.
\]

Such an index exists because \(2^{-j}\to0\); the moving depth eventually
contains it. Given an error exponent \(D\), choose the analytic
precision \(A\) and the localization exponent \(K_0\) with

\[
D+2+C_E+C_F\le A,\qquad 2A+2C_W+344\le K_0. \tag{C5}
\]

Here \(C_W,C_E,C_F\) are the fixed polynomial exponents respectively
bounding the marked weight, the selected-factor coefficients, and the
Fourier integrals of the ordinary-slot cutoffs; their values are
specified in \protect\hyperlink{bilinear}{the bilinear reduction}.
Choose \(A\) first, then \(K_0\).

The derivation in \protect\hyperlink{bilinear}{Section~\ref{bilinear}}
expands the ordinary slots, allocates a selected divisor
\(e\in[EL^{-88},E]\) with \(E=UL^{-K_0}\), and applies Cauchy--Schwarz.
For the explicitly defined sum \(\mathcal B'\) in (B8), with
\(G=L^{C_E}\prod_i|\widehat f_{\delta_i}(\vartheta_i)|\), it gives

\[
|\mathcal B'|^2\le6EVG^2
 \left(2xL^{2C_W+344}+(C_{\rm phys}+24C_1)xL^{-A}\right). \tag{C6}
\]

The first term is the diagonal divisor-moment bound. The second follows
from the \protect\hyperlink{marked-correlation}{marked correlation} and
its Fourier tail. Since \(EV\le C_{\rm box}xL^{-K_0}\), (C5) implies
\(|\mathcal B'|\le CGxL^{-A}\). The remaining slot integrals cost
\(L^{C_E+C_F}\). Fourier inversion of \(\psi(t)=\Psi(e^t)\) restores the
localization; its zeroth and \(H\)th absolute Fourier moments are
finite, so the resulting constant is independent of \(L\) and of the
bounded companion coefficient.

Apply the unsigned and signed estimates to the exact projection
\(A_\rho=(A_0+\rho A_0E)/2\). The triangle inequality and \(|\rho|=1\)
give the same bound for \(\mathcal T_\rho\).

\Needspace{6\baselineskip}

\hypertarget{correlation-two-successive-replacements}{%
\subsection{Two successive
replacements}\label{correlation-two-successive-replacements}}

Write \(\mathcal T_{ij}\) for (C4) with the first coefficient sharp when
\(i=0\) and approximated when \(i=1\), and similarly for the second
coefficient. Then

\[
\mathcal T_{00}-\mathcal T_{11}
=(\mathcal T_{00}-\mathcal T_{10})
 +(\mathcal T_{10}-\mathcal T_{11}). \tag{C7}
\]

On the cells under consideration, (C1) gives \(0\le B_I\le1\). Both
companions are therefore bounded, so (B3) applies in either order. We
use bounds \(Q\) and \(2Q\), respectively, for the two differences.

A rectangle is partitioned into its attained binary dyadic rectangles.
The finite carrier consists of the scales of actual positive factor
pairs \(mn=u+2\) with \(\Psi(u/x)\ne0\). There are at most

\[
D_Y=(\lfloor\log_2(Y+2)\rfloor+1)^2 \tag{C8}
\]

such pairs of scales. All full-cell ratios remain those of (C1). The
internal subdivision cost is absorbed by the extra \(2\) in (C5). The
external rectangle count (C8) is then paid explicitly. In the special
balanced positive estimate, a sharper linear count is available; that
separate count is used in
\protect\hyperlink{sieve-counting-balanced-rectangles}{Section~\ref{sieve-counting-balanced-rectangles}}.

\Needspace{6\baselineskip}

\hypertarget{correlation-the-finite-error-sum-and-its-quantifiers}{%
\subsection{The finite error sum and its
quantifiers}\label{correlation-the-finite-error-sum-and-its-quantifiers}}

Fix an integer \(N\) and a nondecreasing sequence

\[
t_0=b_1,\qquad t_N=1/2-\kappa.
\]

For \(i<N\), the unbalanced region consists of positive pairs \(mn=u+2\)
with \(x^{t_i}<n\le x^{t_{i+1}}\). Its first sharp coefficient is
\(R_{x^{t_i}}(m)\) and its second is the prime indicator at \(n\). For
the balanced region both coefficients are prime indicators, and both
factors lie in \([x^{1/2-2\kappa},x^{1/2+2\kappa}]\).

Let \(U^{(i)}_{ab,\rho}\) and \(B_{ab,\rho}\) be the corresponding sums
with replacement flags \(a,b\in\{0,1\}\). The total error actually used
is

\[
\mathcal E(L)=\sum_{\rho=\pm1}\left(
 \sum_{i<N}|U^{(i)}_{00,\rho}-U^{(i)}_{10,\rho}|
 +|B_{00,\rho}-B_{10,\rho}|
 +|B_{10,\rho}-B_{11,\rho}|\right). \tag{C9}
\]

A second quantity \(\mathcal E_{\rm all}\) sums all four successive
differences, in both replacement orders, for each of the \(N+1\) regions
and both signs.

The bounds have the following quantifier order:

\[
\begin{gathered}
\text{fix }N,b_1,\kappa,D\text{ and the cutoff }\Psi;\\
\text{choose }H,J,B,S,C\text{ with }S>2(B+1)+3, C>0;\\
\text{then for every admissible sequence }(t_i),
\text{ eventually in }L,\\
\mathcal E(L)\le2(N+3)Q(L),\qquad
\mathcal E_{\rm all}(L)\le12(N+1)Q(L),\\
Q(L)=C D_Y xL^{-D}. \tag{C10}
\end{gathered}
\]

The threshold may depend on the partition. The approximation parameters
and \(C\) precede its endpoints. To check the factors in (C10), for one
sign there are \(N\) unbalanced first steps and two balanced steps with
bounds \(Q\) and \(2Q\). This gives \((N+3)Q\). For each region and sign
the four steps in both orders cost \(Q+2Q+Q+2Q=6Q\).

The conclusion follows from the rectangle estimates by finite partition
and the triangle inequality.

\Needspace{6\baselineskip}

\hypertarget{correlation-normalization-by-the-actual-mass}{%
\subsection{Normalization by the actual
mass}\label{correlation-normalization-by-the-actual-mass}}

Take \(D=230\). Since \(Y=\lfloor2e^L\rfloor\), (C8) gives
\(D_Y=O(L^2)\). For fixed \(N\) and the fixed parameters in (C10),

\[
\mathcal E(L),\mathcal E_{\rm all}(L)
 =O(e^LL^{-228})=o(e^LL^{-222}). \tag{C11}
\]

The lower bound from \protect\hyperlink{weights}{Section~\ref{weights}}
is uniform in the sign:

\[
\frac{X_\rho}{L}\ge c_0e^LL^{-222}.
\]

Given \(\varepsilon>0\), apply (C11) with the absolute error tolerance
\(\varepsilon c_0\) and intersect its eventual range with the mass lower
bound. This proves, for either sign,

\[
\mathcal E(L)=o(X_\rho/L),\qquad
\mathcal E_{\rm all}(L)=o(X_\rho/L). \tag{C12}
\]

In particular, (C7) gives the balanced two-step error required in the
final sieve argument.
